\section{Results}
\label{sec:results}

The final comparison contains two variants for each model: without lag 
and with lag features. This design directly answers the central question 
of our project: is the performance gain caused by a more complex 
algorithm, or by better temporal information?

\begin{table}[htbp]
\caption{Final test results after hyperparameter tuning.}
\label{tab:results}
\centering
\scriptsize
\begin{tabularx}{\columnwidth}{p{0.27\columnwidth} l r r r r}
\toprule
\textbf{Model} & \textbf{Lag} & \textbf{MAE} & \textbf{RMSE} & 
\textbf{$R^2$} & \textbf{MAPE} \\
\midrule
Dummy          & no  & 38.14 & 48.02 & -0.499 & 0.817 \\
Linear         & no  & 27.67 & 35.75 &  0.169 & 0.613 \\
Linear         & yes & 19.04 & 24.97 &  0.595 & 0.340 \\
Ridge          & yes & 19.00 & 24.92 &  0.596 & 0.345 \\
Lasso          & yes & 17.78 & 23.86 &  0.630 & 0.291 \\
Decision Tree  & yes & 19.47 & 26.68 &  0.537 & 0.282 \\
KNN            & yes & 19.55 & 25.37 &  0.582 & 0.282 \\
Random Forest  & yes & 17.17 & 23.02 &  0.656 & 0.295 \\
Gradient Boost & yes & \textbf{16.92} & \textbf{22.78} & 
                       \textbf{0.663} & \textbf{0.253} \\
Neural Network & yes & 18.89 & 24.31 &  0.616 & 0.299 \\
\bottomrule
\end{tabularx}
\end{table}

The three best models are Gradient Boosting (RMSE 22.78, $R^2$ 0.663), 
Random Forest (RMSE 23.02, $R^2$ 0.656), and Lasso (RMSE 23.86, 
$R^2$ 0.630), all using lag features. Gradient Boosting achieves the 
best accuracy but requires the longest training time (9.6 seconds). 
Random Forest is nearly as accurate and trains much faster. Lasso is 
the strongest linear model and serves as a transparent benchmark 
against the ensemble methods.

\subsection{Effect of Lag Features}

\begin{table}[htbp]
\caption{Effect of lag features on test RMSE (lower is better).}
\label{tab:lag_effect}
\centering
\scriptsize
\begin{tabularx}{\columnwidth}{p{0.30\columnwidth} r r r}
\toprule
\textbf{Model} & \textbf{No lag} & \textbf{Lag} & 
\textbf{Improvement} \\
\midrule
Linear         & 35.75 & 24.97 & 30.2\% \\
Ridge          & 34.87 & 24.92 & 28.5\% \\
Lasso          & 32.91 & 23.86 & 27.5\% \\
Random Forest  & 29.02 & 23.02 & 20.7\% \\
Gradient Boost & 28.24 & 22.78 & 19.3\% \\
Decision Tree  & 30.58 & 26.68 & 12.8\% \\
KNN            & 27.86 & 25.37 &  9.0\% \\
Neural Network & 26.26 & 24.31 &  7.4\% \\
\bottomrule
\end{tabularx}
\end{table}

Lag features improve every model without exception. The improvement is 
largest for linear models (up to 30.2\%) because the lag variables add 
temporal information that weather and calendar features alone cannot 
capture. Ensemble trees benefit less in relative terms because they can 
already extract some temporal structure through nonlinear splits, but 
they still improve by around 20\%.

\begin{figure}[htbp]
\centering
\includegraphics[width=\linewidth]{test_rmse_lag_delta.png}
\caption{Test RMSE decreases for every model when station-specific lag 
features are added. This supports the decision to prioritize temporal 
representation over adding more algorithms.}
\label{fig:lagdelta}
\end{figure}

\subsection{Hyperparameter Tuning Effect}

\begin{table}[htbp]
\caption{Tuning effect: sklearn defaults vs.\ tuned parameters 
(with lag, test set).}
\label{tab:tuning}
\centering
\scriptsize
\begin{tabularx}{\columnwidth}{p{0.32\columnwidth} r r r r}
\toprule
\textbf{Model} & \textbf{RMSE$_{\text{default}}$} & 
\textbf{RMSE$_{\text{tuned}}$} & \textbf{$R^2_{\text{default}}$} & 
\textbf{$R^2_{\text{tuned}}$} \\
\midrule
Gradient Boosting & 23.50 & 22.78 & 0.641 & 0.663 \\
Neural Network    & 24.99 & 24.31 & 0.594 & 0.616 \\
Lasso             & 23.89 & 23.86 & 0.629 & 0.630 \\
\bottomrule
\end{tabularx}
\end{table}

The tuning results show three different patterns. Gradient Boosting 
improved the most: reducing the learning rate from 0.1 to 0.03 and 
increasing the number of trees from 100 to 400 lowered the test RMSE 
by 0.72 points (-3.1\%). This confirms that slower, more conservative 
learning generalizes better for this dataset. The Neural Network also 
improved clearly (-0.68, -2.7\%), but for a different reason: the 
architecture stayed the same (one hidden layer with 128 neurons), 
while the L2 regularization penalty was increased from 0.0001 to 0.1. 
Forcing the network to learn smaller weights reduced overfitting more 
effectively than adding complexity. Lasso barely changed (-0.03, 
-0.1\%), which shows that the sklearn default of $\alpha = 1.0$ was 
already very close to the optimal value for this dataset.

The overall finding is consistent across all three models: the main 
gain came from better regularization and more conservative learning, 
not from increasing model capacity.

\subsection{Result Interpretation}

Gradient Boosting with lag features is the best final model, reaching 
a test RMSE of 22.78, MAE of 16.92, and $R^2$ of 0.663. This means 
that on average, the model's daily rental prediction for a station 
deviates by about 23 rentals from the true value on unseen test data. 
Random Forest is a strong and faster alternative with nearly identical 
accuracy. Lasso shows that a simple regularized linear model becomes 
competitive once the dataset contains the right temporal features.

The most important result is not the winning model alone. Lag features 
improve every model, and better regularization improves tuning results 
across all three algorithm families. Both findings support the same 
conclusion: for this dataset, improving the input representation and 
controlling model complexity matters more than choosing a more 
sophisticated algorithm.